\section{Introduction}
\label{sec:intro}

Large vision-language models connect a visual encoder to a language decoder and extend instruction following to images and video. A standard image already yields hundreds of visual tokens, and high-resolution crops or video frames make the prefix longer. Standard self-attention grows quadratically with that length, and the cache grows with it. Studies of visual redundancy show that a compact subset can still support many understanding tasks, which makes training-free token pruning a direct route to lower inference cost.

Current pruners divide into three families, and we reveal three coupled limitations.

Salience and interaction scores keep high-scoring locations. FastV uses shallow language-model attention, SparseVLM scores visual tokens with image-related text and recycles information across layers, and FitPrune fits the attention pattern under a compute budget~\cite{fastv,sparsevlm,fitprune}. An individual score does not state whether two kept tokens stand in for each other, so high-score regions still contain a large amount of similar content. Visual-token salience is also depth-dependent. The attention a patch receives, and the representative role it plays, change from a middle layer to a deep layer. A selection that reads only one depth therefore misses complementary information emphasized at the other depth.

Diversity and coverage use feature differences or a representative relation. DART down-weights duplicates of a pivot, DivPrune casts selection as a max-min diversity problem, and SCoRe unifies salience and coverage as a weighted $k$-center~\cite{dart,divprune,score}. These criteria improve image-level retention. They are not fully consistent with the evidence the question needs. Background or edge locations that are weakly related to the question still enter the kept set through feature difference or coverage contribution, and they occupy limited interaction capacity. The local text, attributes, and relations required by the answer are not kept first. The representative set therefore has to be adapted to the question, while the loss of the existing coverage relation stays under an explicit bound.

One-shot compression before the language model fixes the later visual length. VisionZip mixes dominant tokens with context, CDPruner puts instruction relevance into conditional diversity, and CoIn combines salience, cross-modal alignment, and coverage~\cite{visionzip,cdpruner,coin}. After that single cut, every later layer sees the same visual capacity. A small entry discards candidates before cross-modal interaction has formed. A large entry carries the extra tokens into deep layers, where they add little further accuracy and block additional compression. PyramidDrop, SparseVLM, and ViTCoP already shorten the sequence by stages~\cite{pyramiddrop,vitcop}. For a nested schedule, a location that never enters the first set cannot be restored by a later ranking. Entry identity and later capacity have to be chosen together.

We cast visual-token pruning as query-adaptive subset selection under a layered compute budget and a coverage-fidelity constraint. Near a fixed-capacity visual basis, the objective raises question relevance. The constraints limit total coverage distortion and the number of identity edits. Nested sets then contract the visual capacity across interaction stages. This statement links who is admitted to the decoder with how computation is allocated afterward, and it keeps representativeness, task relevance, and inference cost in one problem.

Q-SCOPE is a training-free instance of that problem. Dual-depth salience and descriptors build a regional representative basis. Query-guided exchange (QGE) then replaces a bounded number of identities at constant capacity and rejects an edit whose coverage loss exceeds the allowed distortion. In the decoder, question-to-visual attention (Q-Attn), the attention from question tokens to the current visual positions, rescores the survivors at scheduled cuts. The forward pass still gathers the original visual rows at the selected indices.

On LLaVA-1.5-7B, the completed early schedule removes 94.4\% of the visual token-layer budget and retains 95.66\% of the unpruned six-task relative mean~\cite{llava15}. The same-budget comparison with other methods is reserved in Table~\ref{tab:main} and is not filled from numbers measured under a different protocol.

Our contributions are threefold. First, we reveal three coupled limitations of visual token pruning: depth-dependent salience, the mismatch between visual coverage and question evidence, and a fixed capacity across interaction stages. Second, we formulate coverage-constrained query-adaptive pruning and instantiate it with dual-depth bases, regional quotas, QGE, and Q-Attn refinement. Third, we report the completed six-task result, the exchange trade-off, and the effect of when the first cut occurs, and we leave unmeasured cells blank.
